\documentclass[12pt]{article}
\usepackage[margin=0.75in]{geometry}
\usepackage{lmodern}
\usepackage{fancyhdr}
\usepackage{titlesec}
\usepackage{enumitem}
\usepackage{xcolor}
\usepackage[hidelinks]{hyperref}

\setlength{\parindent}{1.1em}
\setlength{\parskip}{0.35em}
\setlength{\headheight}{15pt}
\titleformat{\section}{\large\bfseries\scshape}{}{0pt}{}
\titleformat{\subsection}{\normalsize\bfseries}{}{0pt}{}
\pagestyle{fancy}
\fancyhf{}
\fancyfoot[C]{\thepage}
\renewcommand{\headrulewidth}{0.4pt}

\fancyhead[L]{Whately Elements of Logic}
\fancyhead[R]{Teacher Guide}
\begin{document}
\begin{center}
{\LARGE\bfseries Teacher Guide: Week 3}\par\vspace{0.25em}
{\large\scshape What We Say and What We Mean: Making Clear Claims}\par\vspace{0.4em}
{12th Grade Long Logic Unit Based on Richard Whately's \textit{Elements of Logic}}\par\vspace{0.6em}
\rule{0.82\textwidth}{0.4pt}
\end{center}

\section*{Week 3 Overview}
Week 3 moves from terms (Week 2) to propositions—the statements we make when we combine terms. Students learn to identify the subject, predicate, quality, and quantity of claims, and to translate ordinary language into the four standard forms (A, E, I, O). The week emphasizes precision in stating what is being claimed about \emph{how much} of a subject and \emph{in what way} (affirmative or negative). 

The new test model uses four category "riddles" (Mammals, Birds, Insects, Planets) in which students must spot and repair overreaching claims that misuse "All," "None," or "Most." This directly practices distinguishing universal from particular quantity and recognizing when scope errors make a proposition false even if the subject-predicate link is plausible. The revised twelfth-grade plan also asks students to translate claims from \textit{Macbeth} and \textit{Julius Caesar}, since literary interpretation often turns on imprecise claims such as “Macbeth is fated” or “Caesar is ambitious.” This foundation is essential for the study of syllogisms, mood, and figure in later weeks.

\section*{Essential Question}
How can we state claims clearly enough that others know exactly what we are asserting—especially how much of the subject we are talking about?

\section*{Student-Facing Materials}
\begin{itemize}[leftmargin=*]
\item \textit{What We Say and What We Mean: Making Clear Claims} — student reading handout.
\item \textit{Week 3 Argument Lab: Identifying and Clarifying Propositions} — practice in identifying quality, quantity, distribution, and translating ordinary language.
\item \textit{Macbeth Lit-Anchor: What Does ``A Man'' Mean?} — Shakespeare anchor passage for translating dramatic claims about manhood, courage, murder, fear, and moral limits into precise proposition form.
\item \textit{Week 3 Logic Test: The Professor's Riddles} — assessment using four category lists to identify accurate vs. overreaching AEIO claims.
\end{itemize}

\section*{Learning Goals}
Students will be able to:
\begin{enumerate}[leftmargin=*]
\item Define a proposition and distinguish it from questions, commands, and exclamations.
\item Identify the subject, predicate, and copula of a proposition.
\item Distinguish affirmative from negative propositions (quality).
\item Distinguish universal from particular propositions (quantity).
\item Recognize the four standard forms (A, E, I, O).
\item Determine whether terms are distributed or undistributed.
\item Translate common English expressions into standard logical form.
\item Spot and repair overreaches in quantity ("All" when only "Some," "None" when exceptions exist, "Most" when unsupported).
\item Correctly state the distribution pattern for O propositions: the subject is undistributed and the predicate is distributed.
\item Translate recurring Shakespeare claims into precise proposition form for the course portfolio.
\end{enumerate}

\section*{Key Vocabulary}
proposition, subject, predicate, copula, quality, quantity, universal, particular, distribution, affirmative, negative, A/E/I/O forms, scope / overreach

\section*{Weekly Lesson Sequence}

\subsection*{Day 1: What Is a Proposition?}
\begin{itemize}[leftmargin=*]
\item Begin with examples of statements that can be true or false versus questions and commands.
\item Introduce the idea that logic is concerned with claims that can be evaluated.
\item Read the first half of the student reading.
\item Students identify five propositions from everyday language and explain why they qualify.
\end{itemize}

\subsection*{Day 2: Quality and Quantity}
\begin{itemize}[leftmargin=*]
\item Focus on the distinction between affirmative/negative and universal/particular.
\item Use multiple examples on the board, including simple "All," "Some," "No," and "Most" claims. Have students classify each one.
\item Complete the first two sections of the Argument Lab in pairs. Introduce the idea that "Most" is a particular claim that can still be too strong.
\end{itemize}

\subsection*{Day 3: The Four Forms and Distribution}
\begin{itemize}[leftmargin=*]
\item Introduce the A, E, I, O labels with examples.
\item Practice determining distribution of terms.
\item Students work through Part 2 of the Lab and explain their reasoning.
\item Preview the test riddles: show one short list and model spotting a false "All" vs. true "Some."
\end{itemize}

\subsection*{Day 4: Translating Ordinary Language and Spotting Scope Errors}
\begin{itemize}[leftmargin=*]
\item This is often the most difficult skill. Spend significant time on translation.
\item Work through Part 3 as a class, then have students complete Part 4 independently.
\item Use the Macbeth lit-anchor to show how the ordinary-language challenge ``be a man'' can hide a universal proposition that needs to be stated and tested.
\item Introduce scope critique: practice rewriting an overreaching claim from the riddles into correct AEIO form.
\item Discuss why translation and accurate quantity matter for clear argument and avoiding false generalizations.
\end{itemize}

\subsection*{Day 5: Window Notes, Riddles Practice, and Review}
\begin{itemize}[leftmargin=*]
\item Students complete the portfolio window note task from the Lab, preferably using one proposition from \textit{Macbeth} or \textit{Julius Caesar}. If using the Macbeth lit-anchor, have them complete the Macbeth proposition translation chart.
\item Class activity or homework: students work a shortened version of one or two riddles from the test, marking falses and repairing one.
\item Class discussion: How does making quantity and quality explicit change the way an argument can be supported or challenged? Why do overreaches with "All" or "None" matter in science and everyday reasoning?
\item Exit task: Write one claim from the unit in two different logical forms (e.g., change All to Some) and explain the difference in what it asserts.
\end{itemize}

\section*{Lit-Anchor Teaching Notes}
The Week 3 lit-anchor focuses the proposition lesson on Lady Macbeth's pressure over manhood. Do not let the class remain at the level of insult or gender-role discussion alone. The logic task is to translate dramatic language into precise propositions before testing it.

\begin{itemize}[leftmargin=*]
\item \textbf{Guided passage:} \textit{Macbeth} Act 1, Scene 7, Lady Macbeth's manhood challenge and Macbeth's reply: ``I dare do all that may become a man; / Who dares do more is none.''
\item \textbf{Independent passage:} Act 3, Scene 4, the banquet scene where Lady Macbeth asks, ``Are you a man?'' and Macbeth answers, ``What man dare, I dare.''
\item \textbf{Central logical pressure:} a vague challenge like ``be a man'' hides propositions about courage, murder, fear, and moral limits.
\item \textbf{Portfolio artifact:} Macbeth proposition translation chart: original wording, ordinary-language claim, standard logical form, subject, predicate, quality, quantity, and what becomes clearer.
\end{itemize}

Strong guided answers should identify Lady Macbeth's implied claim as something like: ``All true men dare this deed,'' or ``All truly manly people are people who dare the murder.'' That is a universal affirmative claim (A). Macbeth's reply can be stated as a universal negative: ``No people who dare beyond what befits a man are true men.'' (E) A useful counter-proposition is: ``Some daring acts are not honorable acts.'' (O) This O proposition helps students separate courage from goodness and shows why clear proposition form matters.

When students reach the banquet passage, look for the change in Macbeth. Earlier he briefly defines manhood by moral limit; later he speaks more in terms of boldness, fearlessness, and composure. That contrast gives students a literary reason to care about proposition form: the same word can carry a changed predicate across the play.

\section*{Notes on the Logic Test (New Model)}
The test replaces the previous long translation story with four compact "riddles"—lists of ten statements each about a scientific category. The framing is the Department of Science professor asking students to help locate where her riddles go wrong. The core task practices exactly the week's focus: accurate use of quantity (universal vs. particular) in AEIO propositions.

\textbf{Test structure:}
- Brief story and Whately focusbox on quantity/quality.
- Part I: The Four Riddles (Mammals, Birds, Insects, Planets). Students mark T/F for each statement's accuracy as a description of the category and record false numbers with AEIO explanations.
- Spiral Parts II–III: Integrate Week 1 (what follows) and Week 2 (shifting terms) using the riddle statements.
- Part IV: Repair one false claim (rewrite with correct quantifier + supporting fact).
- Part V: What logic can test vs. needs external facts (distribution, quality change, scope vs. discovery).
- Final reflection on why precise "all/some/none/most" matters.

\textbf{Answers for Part I (False statements and key teaching notes):}

\textbf{Riddle 1 — Mammals}\\
False: 3, 5\\
- 3: "None can live in water" (E claim) is false. Some mammals (whales, dolphins, otters) live in water. Correct: Some (I) can live in water.\\
- 5: "All give birth to live young" (A claim) is false. Monotremes (platypus, echidna) lay eggs. Correct: Some (I) give birth to live young; or Most (strong particular) do.\\
Debatable/tricky: 4 "Most have hair or fur..." — generally true but whales have very little; excellent discussion of when "most" is safe vs. overclaim. 2 "Some lay eggs" is true (monotremes).

\textbf{Riddle 2 — Birds}\\
False: 3, 8, 10\\
- 3: "None can swim" (E) false. Penguins, ducks, etc. swim. Some (I) can swim.\\
- 8: "All have hollow bones" (A) is tricky/false in strict sense. Many birds do, but not literally every species or every bone in a simple way; not a universal biological rule without qualification.\\
- 10: "None are warm-blooded" (E) false. Birds are endothermic (warm-blooded). All (A) or Most are.\\
Note: 2, 5, 7, 9 are classic "Some" or "Most" truths (flightless birds, migration). 6 "None have teeth as adults" is true for modern birds.

\textbf{Riddle 3 — Insects}\\
False: 4, 6, 8\\
- 4: "None live in water" (E) false. Many aquatic insects (dragonfly larvae, water striders, etc.). Some (I) live in water.\\
- 6: "All undergo metamorphosis" (A) is especially tricky. Many do (complete or incomplete), but some have direct development or minimal change. Not strictly All.\\
- 8: "None have antennae" (E) false. Insects characteristically have antennae. All (A) have them.\\
Note: 1, 2, 10 are solid universals for the category definition. 7 colonies (ants, bees) is good "Some."

\textbf{Riddle 4 — Planets}\\
False: 4, 7, 10\\
- 4: "All are made of rock" (A) false. Gas giants (Jupiter, Saturn) and ice giants are not primarily rock. Some (terrestrial) are; others are not.\\
- 7: "None have atmospheres" (E) false. Earth, Mars, Venus, gas giants all have atmospheres (varying). Some (I) or Most do.\\
- 10: "None can support storms" (E) false. Jupiter's Great Red Spot, etc. Some (I) have massive storms.\\
Tricky: 6 "Most spin on an axis" — true enough in practice; 8 "Some are larger than Earth" true (gas giants).

\textbf{Teaching emphasis:} The mistakes are usually in \textbf{scope}:
\begin{itemize}
\item \textbf{All} when it should be \textit{most} or \textit{some} (universal overclaim).
\item \textbf{None} when exceptions exist (universal negative too strong).
\item \textbf{Most} when evidence is unclear or the claim is still too broad for the context.
\end{itemize}
That is exactly the habit the week (and later syllogism work) wants students to catch: check the quantifier before accepting or building on a claim.

The test also reinforces that logic can expose the \textit{form} of the overreach (wrong quantity) even when the underlying facts require outside knowledge to confirm which statements are actually false.

\textbf{Grading suggestion:} Credit for correctly identifying the falses (with some tolerance on debatable items like Mammals 4 or Birds 8) plus quality of AEIO explanations and repairs. Spiral parts check retention.

\section*{Notes on Examples}
The student reading deliberately provides numerous examples of each concept. Students benefit from seeing the same idea illustrated in several different contexts. The translation exercises in the Lab are especially important; many students initially struggle to move between ordinary language and standard form. Encourage them to ask “What exactly is being claimed about every member, some members, or no members of the subject?”

The riddles in the test extend this by requiring students to apply the same question to real scientific categories and defend their judgments about scope.

\section*{Connection to Future Units}
This week prepares students for syllogistic reasoning. In the next week, students will learn to evaluate whether a conclusion follows from two premises. That evaluation depends entirely on the ability to identify the quality, quantity, and distribution of the propositions involved. Accurate scope (avoiding hidden All/None errors) becomes even more critical when two propositions are chained. The following week should then add mood and figure after students have practiced the practical structure of syllogisms. The riddle exercise directly foreshadows spotting undistributed terms or illicit process later.

\end{document}